\documentclass[10pt,a4paper]{amsart}
\usepackage[margin=19mm]{geometry}
\usepackage{amsmath,amssymb,mathtools}
\usepackage[scr=rsfs,cal=euler]{mathalfa}
\usepackage{xfrac}
\usepackage[unicode,hidelinks,bookmarks=false]{hyperref}
\newcommand{\ie}{i.e.\ }
\newcommand{\cf}{cf.\ }
\newcommand{\resp}{respectively\ }
\newcommand{\Subsection}{\subsection}
\providecommand{\nobreakdash}{\nobreak\hbox{-}}
\setlength{\emergencystretch}{2em}
\multlinegap0pt
\newtheorem{proposition}{Proposition}
\newcommand{\Z}{\mathbb{Z}}
\title[completeness proposition]{On the heat kernel of a Cayley graph of $\operatorname{PSL}_2\mathbb{Z}$: the completeness of the eigenfunction family}
\author{Anders Karlsson, Kamila Kashaeva}
\date{Source: arXiv:2506.02340v3 (CC BY 4.0)}
\begin{document}
\maketitle

\noindent\emph{Source and licence.} On the heat kernel of a Cayley graph of $\operatorname{PSL}_2\mathbb{Z}$, by Anders Karlsson, Kamila Kashaeva. Version arXiv:2506.02340v3
(CC BY 4.0), distributed by arXiv under the Creative Commons Attribution 4.0 International licence,
http://creativecommons.org/licenses/by/4.0/, as recorded in the arXiv licence metadata for this exact
version and shown on the abstract page https://arxiv.org/abs/2506.02340v3. Full licence notice:
\texttt{licenses/arxiv-2506.02340-CC-BY-4.0.txt}.\par\smallskip

\noindent\textit{Editorial scope.} Counted unit: the article's completeness proposition (source0.tex lines 1524--1529, source label \texttt{completeness1}): for all integers $m,n$ the discrete eigenfunctions together with the continuous family of Eisenstein-type functions satisfy the displayed orthonormality relation whose right-hand side is the Kronecker delta. The article's complete original proof of it is delivered with it (source0.tex lines 1532--1613, one \texttt{proof} environment of 82 lines, adjacent to the statement, with no gap and no deferral). No mathematics is written, repaired or re-derived: the statement and the proof are the article's own lines, in the article's own notation, and no retained line is substituted, re-flowed or omitted, so nothing is removed and the package carries no omission record. Every cross-reference of every retained line resolves inside this document, so no further source-local context is needed; no retained line contains a citation, so the wrapper carries no bibliography and no bibliography entry; and no retained line contains a figure environment, an image-inclusion command or drawing code, so no image is delivered and the package declares no asset dependency.

Editorial formatting only, in addition: an AMS \texttt{amsart} preamble that restates the article's own declarations and macros used by the retained lines, namely the \texttt{proposition} environment and the macro \texttt{\textbackslash Z} copied from the article's own preamble (source0.tex line 18). The retained environment prints in this document as Proposition 1, whereas the article prints the same item with the numbering of its own full document; the numbered displays of the retained lines print consecutively in order of appearance, whereas the article numbers its equations per section.

The complete native source of the article as distributed in the arXiv e-print for this exact version is bundled unchanged as \texttt{sources/179199/source0.tex}: it is byte identical to the e-print file \texttt{heatkernelPSL2Z\_arXiv\_may26.tex} except that the pipeline's mechanical concatenation prepends one header comment line naming it. The article uses the AMS \texttt{amsart} class; no publisher class or style file is shipped in the e-print and none is redistributed or used.
\begin{proposition}\label{completeness1}  
The following completeness condition is satisfied: for all $m,n\in\Z$,
$$
\sum_{\epsilon\in\{\pm1\}}f_\epsilon(m)f_\epsilon(n)+\int_0^\pi \sum_{\mu,\epsilon\in\{\pm1\}}f_{x,\mu,\epsilon}(m)f_{x,\mu,\epsilon}(n)\frac{\mathrm{d}x}{H_{\mu\epsilon}(x)}=\delta_{m,n}.
$$
\end{proposition}

\begin{proof}
 Let us denote
\begin{equation}
 F_x(m,n):=\sum_{\mu,\epsilon\in\{\pm1\}}f_{x,\mu,\epsilon}(m)f_{x,\mu,\epsilon}(n)\frac{1}{H_{\mu\epsilon}(x)} \label{F_x}
\end{equation}
 and
\begin{equation}
  C(m,n):=\int_0^\pi F_x(m,n) \ \mathrm{d}x +\sum_{\epsilon\in\{\pm1\}}f_\epsilon(m)f_\epsilon(n). \label{C_x}
\end{equation}
 We need to check separately the different cases depending on the parity of $m$ and $n$. 
The cases where both entries are negative are equivalent to the cases with both positive entries because of the symmetry property of the (generalized) eigenfunctions
$$
f_\epsilon(-m-1)=\epsilon f_\epsilon(m), \quad f_{x,\mu, \epsilon}(-m-1)=\epsilon f_{x,\mu, \epsilon}(m), \quad \forall m\in\Z.
$$
We will show in detail only the computation of $C(2m,2n+1)$ with $m,n\in\Z_{\ge0}$ since 
the computations of the other cases are analogous.
First, we compute the sum in \eqref{F_x} with the result
\begin{equation}\label{eq:positive-even-odd}
F_x(2m,2n+1)=\frac{\sqrt{2}}{\pi} \cdot \frac{2\big( \sin(mx)\sin(nx)-\sin((m+1)x)\sin((n+1)x)\big)}{4\cos(2x)-5}.
\end{equation}
Next, we transform  the numerator and denominator of~\eqref{eq:positive-even-odd} separately. The identity for the product of sines
$$
2\sin(x)\sin(y)=\cos(x-y)-\cos(x+y)
$$
allows us to rewrite the numerator of~\eqref{eq:positive-even-odd}  as
\begin{align}
 2\big( \sin(mx)\sin(nx) & -\sin((m+1)x)\sin((n+1)x)\big) = \cos(x(m+n+2))-\cos(x(m+n)) \notag \\
 & = \operatorname{Re}(e^{ix(m+n+2)}-e^{ix(m+n)}) = \operatorname{Re}(e^{ix(k+2)}-e^{ixk}), \label{num_ee}
\end{align}
where $k:=m+n$.
Denoting $\xi:=e^{2ix}$, we rewrite the denominator of~\eqref{eq:positive-even-odd}  as
$$
4\cos(2x)-5=2(\xi+\xi^{-1})-5=\frac{2}{\xi}(\xi-2)(\xi-\frac{1}{2}).
$$
Using the fraction decomposition for the inverse and the geometric series (expanding for $|\xi|=1$), we obtain
\begin{align}
 \frac{1}{4\cos(2x)-5} &=\frac{\xi/2}{(\xi-2)(\xi-\frac{1}{2})} 
=-\frac{1}{3}\Big(\frac{1}{1-\frac{\xi}{2}}+\frac{1}{2\xi}\frac{1}{1-\frac{1}{2\xi}}\Big) =-\frac{1}{3}\Big(\sum_{l=0}^{\infty}\big(\frac{\xi}{2}\big)^l + \sum_{l=0}^{\infty} \big(\frac{1}{2\xi}\big)^{l+1} \Big) \notag \\ 
&=-\frac{1}{3}\Big(\sum_{l=0}^{\infty}2^{-l}\xi^l + \sum_{l=-\infty}^{-1} (2\xi)^l \Big) 
=-\frac{1}{3} \sum_{l\in\Z}2^{-|l|}\xi^l=-\frac{1}{3} \sum_{l\in\Z}2^{-|l|}e^{2ixl}. \label{denom_ee}
\end{align}
Using \eqref{num_ee} and \eqref{denom_ee}, we can rewrite $F_x$ as follows:
\begin{align*}
 F_x(2m,2n+1)
 &=\frac{\sqrt{2}}{\pi} \operatorname{Re}\Big(\frac{1}{4\cos(2x)-5}(e^{ix(k+2)}-e^{ixk})\Big) \\
 &=-\frac{\sqrt{2}}{3\pi} \sum_{l\in\Z}2^{-|l|} \operatorname{Re}(e^{ix(2l+k+2)}-e^{ix(2l+k)}) \\
 &=-\frac{\sqrt{2}}{3\pi} \sum_{l\in\Z}2^{-|l|} \big( \cos(x(2l+k+2))-\cos(x(2l+k)) \big).
\end{align*}
Using the fact that 
$$
\int_0^\pi \cos(x(2l+k)) \mathrm{d}x =
\begin{cases}
\pi \delta_{l,-k/2} & k \ \text{even} \\
0 & \text{else},
\end{cases}
$$
we can now compute the integral (recall that $k=m+n$)
\begin{align*}
 \int_0^\pi F_x(2m,2n+1) \mathrm{d}x & =-\frac{\sqrt{2}}{3\pi} \sum_{l\in\Z}2^{-|l|} \int_0^\pi \cos(x(2l+k+2))-\cos(x(2l+k)) \mathrm{d}x \\
 &= -\frac{\sqrt{2}}{3\pi}
\begin{cases}
 \pi 2^{-(k+2)/2} - \pi 2^{-k/2} & k \ \text{even} \\
 0 & \text{else}
\end{cases}
\\
&= 
\begin{cases}
\frac{\sqrt{2}}{6} 2^{-k/2}  & k \ \text{even} \\
 0 & \text{else}.
 \end{cases}
\end{align*}
Finally, using the definition of the eigenfunctions $f_\epsilon(m)$, we compute the sum in \eqref{C_x}
\begin{align*}
 \sum_{\epsilon\in\{\pm1\}}f_\epsilon(2m)f_\epsilon(2n+1) & =\sum_{\epsilon\in\{\pm1\}}-\frac{1}{6\sqrt{2}}(-\sqrt{2}\epsilon)^{-k}=-\frac{1}{6\sqrt{2}}(-1)^k\sqrt{2}^{-k}(1+(-1)^k) \\
 &=
 \begin{cases}
  -\frac{\sqrt{2}}{6}\sqrt{2}^{-k} & k \ \text{even} \\
 0 & \text{else}.
\end{cases}
\end{align*}
Thus, we have shown that $C(2m,2n+1)= 0$.
\end{proof}

\end{document}
